\Qhead{1.9}{Собственные и вынужденные колебания линейных распределённых систем}{Собственные и вынужденные колебания линейных распределенных систем. Собственные функции системы (моды). Разложение вынужденных колебаний по системе собственных функций.}
\Subfile{1.9.1}{sections/01_oscillations/q09/s01}
\Subfile{1.9.2}{sections/01_oscillations/q09/s02}

\recap{Границы определяют моды, самосопряжённость --- ортогональность, проекция источника --- возбуждение. Модальное разложение сводит вынужденную задачу к осцилляторам.}
\board{Первые моды струны; собственная задача; интеграл проекции силы и сумма откликов.}
\oral{Почему точечная сила на собственной частоте может не возбудить моду?}
\bait{Можно оставить эквивалентность функции Грина и суммы мод: обе дают отклик на тот же источник.}
\sources{Биленко, Воронцов, Вятчанин~\cite{bilenko}, §2.3.1, с.~76--82: распределённая длинная линия; Белокопытов~\cite{belokopytov2}, §1 --- модальная постановка. Обобщение на оператор, нормировка и вынужденное разложение выведены здесь. Отдельный специальный курс Биленко из воспоминания рецензента не идентифицирован.}
